\documentclass[11pt,a4paper]{article}

% Page layout similar to NeurIPS
\usepackage[margin=1in]{geometry}
\usepackage[utf8]{inputenc}
\usepackage[T1]{fontenc}
\usepackage{hyperref}
\usepackage{url}
\usepackage{booktabs}
\usepackage{amsfonts}
\usepackage{nicefrac}
\usepackage{microtype}
\usepackage{amsmath}
\usepackage{amssymb}
\usepackage{graphicx}
\usepackage{xcolor}
\usepackage{algorithm}
\usepackage{algorithmic}
\usepackage{multirow}

\DeclareMathAlphabet{\mathsfit}{T1}{cmss}{m}{sl}
\SetMathAlphabet{\mathsfit}{bold}{T1}{cmss}{b}{sl}
\newcommand{\E}{\mathbb{E}}
\newcommand{\R}{\mathbb{R}}
\newcommand{\cN}{\mathcal{N}}
\newcommand{\cL}{\mathcal{L}}
\newcommand{\cA}{\mathcal{A}}
\newcommand{\cP}{\mathcal{P}}
\newcommand{\given}{\,|\,}

\title{Feature-Diversity-Aware Supervised Contrastive Learning: \\ Mitigating Feature Suppression through Adaptive Pair Weighting}

\author{Anonymous Author\\
Anonymous Institution\\
\texttt{anonymous@example.com}}

\begin{document}

\maketitle

\begin{abstract}
Supervised Contrastive Learning (SCL) has demonstrated remarkable success in learning discriminative representations by pulling together samples from the same class while pushing apart different classes. However, SCL suffers from a critical limitation known as \emph{feature suppression} (or class collapse): the model learns only a subset of discriminative features within each class while suppressing diverse features that could distinguish subclasses or capture intra-class variation. This phenomenon significantly compromises representation quality for fine-grained and downstream tasks. We propose FD-SCL (Feature-Diversity-Aware Supervised Contrastive Learning), a novel method that addresses feature suppression through adaptive pair weighting. Our key insight is that positive pairs sharing rare or underrepresented features should be upweighted to prevent these features from being suppressed by dominant ones. FD-SCL introduces a lightweight Feature Diversity Estimation Module (FDEM) that estimates feature rarity from batch statistics and computes diversity-aware weights without requiring additional labels or auxiliary information. On CIFAR-100 coarse-to-fine evaluation, FD-SCL improves fine-grained accuracy from 1.75\% (standard SCL) to 6.20\%, representing a 254\% relative improvement in feature preservation, while maintaining comparable coarse accuracy. FD-SCL adds less than 5\% computational overhead compared to standard SCL, making it a practical drop-in replacement for existing frameworks.
\end{abstract}

\section{Introduction}

Supervised Contrastive Learning (SCL) \citep{khosla2020supervised} has emerged as a powerful paradigm for representation learning, demonstrating superior performance over traditional cross-entropy training on various benchmarks. By constructing positive pairs from samples sharing the same label and negative pairs from different classes, SCL learns embeddings that are both discriminative and generalizable.

However, a fundamental limitation of SCL has been recently identified: \textbf{feature suppression} (also called class collapse) \citep{xue2023which, zhang2024learning}. When training with SCL, models tend to learn only a subset of discriminative features within each class while suppressing other potentially valuable features. For example, when learning representations of ``dogs,'' the model may focus predominantly on common features like ``four legs'' and ``fur'' while suppressing diverse features like ``tail shape,'' ``ear position,'' or ``coat patterns'' that could distinguish different breeds or poses. This phenomenon significantly compromises representation quality, especially for downstream tasks requiring fine-grained discrimination or handling intra-class diversity.

The root cause of feature suppression lies in the SCL loss formulation: all positive pairs (samples from the same class) are weighted equally. This uniform treatment causes the model to optimize for the ``easiest'' shared features across all positive pairs, effectively averaging out diverse features that are not shared by all samples in the class. As the model converges, representations of subclasses within a class become indistinguishable---a phenomenon known as class collapse \citep{xue2023which}.

\textbf{Key Insight:} Not all positive pairs in SCL are equally informative for learning diverse features. Pairs that share rare or underrepresented features should be upweighted to prevent these features from being suppressed by more dominant ones. By adaptively weighting positive pairs based on their feature diversity contribution, we can learn more comprehensive representations that capture the full spectrum of intra-class variation.

\textbf{Our Contributions:}
\begin{itemize}
    \item We propose FD-SCL, a novel supervised contrastive learning method that addresses feature suppression through adaptive pair weighting based on feature diversity estimation.
    \item We introduce the Feature Diversity Estimation Module (FDEM), which computes diversity-aware weights from batch statistics without requiring additional labels, auxiliary models, or multi-stage training.
    \item We demonstrate that FD-SCL significantly reduces feature suppression: on CIFAR-100 coarse-to-fine evaluation, FD-SCL achieves 6.20\% fine-grained accuracy compared to 1.75\% for standard SCL---a 254\% relative improvement---while maintaining comparable coarse accuracy (83.12\% vs. 83.64\%).
    \item We show that FD-SCL adds less than 5\% computational overhead compared to standard SCL, making it practical for large-scale applications unlike multi-stage alternatives that require 2--3$\times$ compute.
\end{itemize}

\section{Related Work}

\paragraph{Supervised Contrastive Learning.}
\citet{khosla2020supervised} introduced the foundational supervised contrastive loss, demonstrating superior performance over cross-entropy on ImageNet. Subsequent work explored hard negative sampling \citep{robinson2021contrastive, jiang2024hardnegative}, but these focus on selecting difficult negatives rather than addressing feature diversity within positives.

\paragraph{Feature Suppression and Class Collapse.}
\citet{xue2023which} provided the first theoretical analysis of class collapse, showing that SGD's simplicity bias causes models to learn only easy features and suppress hard class-relevant features. \citet{zhang2024learning} proposed Multistage Contrastive Learning (MCL) to address this through progressive feature learning across multiple stages. However, MCL requires training the model 2--3 times sequentially, significantly increasing computational cost.

\paragraph{Temperature Adaptation and Multi-Head Methods.}
Recent work has explored adaptive temperature strategies \citep{kukleva2023temperature, wang2024adaptive} and multi-head architectures for contrastive learning. These methods focus on alignment difficulty rather than feature diversity within classes. Our approach is complementary: FD-SCL can be combined with adaptive temperature methods for further improvements.

\paragraph{Intra-Class Diversity Methods.}
Several methods explicitly encourage intra-class diversity \citep{duboudin2021encouraging, chen2023coinseg}, but these are designed for specific domains (domain generalization, segmentation) or require additional supervision (masks, auxiliary labels). In contrast, FD-SCL is domain-agnostic and requires only standard class labels.

\begin{table}[h]
\caption{Comparison with related methods. FD-SCL addresses feature suppression with minimal overhead.}
\label{tab:comparison}
\centering
\begin{tabular}{lcccc}
\toprule
Method & Addresses & Multi-Stage & Domain- & Overhead \\
& Feature Supp. & Required & Agnostic & \\
\midrule
Standard SCL \citep{khosla2020supervised} & \xmark & \xmark & \cmark & --- \\
MCL \citep{zhang2024learning} & \cmark & \cmark (2--3$\times$) & \cmark & High \\
AMCL \citep{wang2024adaptive} & Partial & \xmark & \cmark & Medium \\
\textbf{FD-SCL (Ours)} & \cmark & \xmark & \cmark & Low ($<$5\%) \\
\bottomrule
\end{tabular}
\end{table}

\section{Method}

\subsection{Preliminaries: Standard Supervised Contrastive Learning}

Given a batch of $N$ samples $\{(x_i, y_i)\}_{i=1}^N$ where $y_i$ is the class label, SCL creates two augmented views for each sample, resulting in $2N$ augmented samples. Let $z_i \in \R^D$ denote the $\ell_2$-normalized embedding of sample $i$ from the projection head. The supervised contrastive loss for anchor $i$ is:
\begin{equation}
\cL_{\text{SCL}}^{(i)} = \frac{-1}{|\cP(i)|} \sum_{p \in \cP(i)} \log \frac{\exp(z_i \cdot z_p / \tau)}{\sum_{a \in \cA(i)} \exp(z_i \cdot z_a / \tau)}
\end{equation}
where $\cP(i) = \{p : y_p = y_i, p \neq i\}$ is the set of positive samples (same class), $\cA(i)$ is the set of all samples excluding $i$, and $\tau$ is the temperature parameter.

\subsection{Feature Diversity Estimation Module (FDEM)}

The core innovation of FD-SCL is estimating feature diversity contribution for each positive pair. We compute this based on \textbf{activation patterns} in the embedding space:

\paragraph{Step 1: Feature Activation Statistics.}
For each dimension $d$ of the normalized embedding $z \in \R^D$, we compute its activation frequency across the batch:
\begin{equation}
f_d = \frac{1}{|B|} \sum_{j \in B} \mathbb{1}[z_{j,d} > \eta]
\end{equation}
where $\eta$ is a threshold (typically 0) and $B$ is the set of samples in the batch.

\paragraph{Step 2: Pairwise Feature Rarity Score.}
For a positive pair $(i, p)$, we compute their shared rare feature score:
\begin{equation}
r_{ip} = \sum_{d=1}^D \frac{\mathbb{1}[(z_{i,d} > \eta) \land (z_{p,d} > \eta)]}{f_d + \epsilon}
\end{equation}
This upweights dimensions that are active in both samples but rare in the batch. The division by $f_d$ ensures that rare features (small $f_d$) contribute more to the rarity score.

\paragraph{Step 3: Diversity-Aware Weight.}
The final weight for positive pair $(i, p)$ is:
\begin{equation}
w_{ip} = \text{softmax}\left(\frac{r_{ip}}{\tau_w}\right)_{p \in \cP(i)} \cdot |\cP(i)|
\end{equation}
where $\tau_w$ is a temperature parameter controlling weight distribution sharpness. The softmax is computed over positive pairs for each anchor, and scaling by $|\cP(i)|$ preserves the original loss scale.

\subsection{FD-SCL Loss}

The FD-SCL loss replaces uniform positive pair weighting with diversity-aware weights:
\begin{equation}
\cL_{\text{FD-SCL}}^{(i)} = \frac{-1}{|\cP(i)|} \sum_{p \in \cP(i)} w_{ip} \cdot \log \frac{\exp(z_i \cdot z_p / \tau)}{\sum_{a \in \cA(i)} \exp(z_i \cdot z_a / \tau)}
\end{equation}

The complete training procedure is summarized in Algorithm~\ref{alg:fd-scl}.

\begin{algorithm}[t]
\caption{FD-SCL Training Procedure}
\label{alg:fd-scl}
\begin{algorithmic}[1]
\REQUIRE Training data $\{(x_i, y_i)\}$, encoder network $f_\theta$, projector $g_\phi$
\REQUIRE Temperature $\tau$, weight temperature $\tau_w$, threshold $\eta$, epochs $T$
\FOR{$t = 1$ to $T$}
    \FOR{mini-batch $\{(x_i, y_i)\}_{i=1}^N$}
        \STATE Create two augmented views: $\{\tilde{x}_i^1, \tilde{x}_i^2\}_{i=1}^N$
        \STATE Compute embeddings: $z_i^k = g_\phi(f_\theta(\tilde{x}_i^k))$
        \STATE Normalize: $z_i^k \leftarrow z_i^k / \|z_i^k\|_2$
        \STATE Compute activation frequencies $f_d$ via Eq.~(2)
        \STATE Compute rarity scores $r_{ip}$ via Eq.~(3)
        \STATE Compute diversity weights $w_{ip}$ via Eq.~(4)
        \STATE Compute FD-SCL loss $\cL_{\text{FD-SCL}}$ via Eq.~(5)
        \STATE Update $\theta, \phi$ via backpropagation
    \ENDFOR
\ENDFOR
\STATE \textbf{return} trained encoder $f_\theta$
\end{algorithmic}
\end{algorithm}

\subsection{Theoretical Analysis}

\textbf{Theorem (Feature Diversity Preservation):} Under the assumption that embeddings follow a mixture distribution where each class contains $K$ subclasses with distinct feature patterns, the expected gradient of FD-SCL preserves gradients for rare features that would be suppressed by standard SCL.

\textbf{Proof Sketch:} Standard SCL gradients for feature dimension $d$ are proportional to $f_d(1-f_d)$. When $f_d \ll 1$ (rare features), gradients vanish. FD-SCL rescales by $1/f_d$, preserving gradient magnitudes for rare features. The full proof is provided in Appendix~A.

\subsection{Computational Complexity}

FDEM adds minimal computational overhead:
\begin{itemize}
    \item Feature activation statistics: $O(B \times D)$
    \item Pairwise rarity scores: $O(|\cP| \times D)$ where $|\cP|$ is the number of positive pairs
    \item Total overhead: $<5\%$ increase in training time compared to standard SCL
\end{itemize}

\section{Experiments}

\subsection{Experimental Setup}

\paragraph{Datasets.}
We evaluate on CIFAR-100 \citep{krizhevsky2009learning}, which provides both coarse (20 superclasses) and fine (100 classes) labels. This enables our primary evaluation: training with coarse labels and measuring fine-grained accuracy to assess feature suppression. We also report results on standard CIFAR-100 classification with 100 fine labels.

\paragraph{Baselines.}
We compare against:
\begin{itemize}
    \item \textbf{Cross-Entropy (CE)}: Standard supervised learning
    \item \textbf{SCL} \citep{khosla2020supervised}: Standard supervised contrastive learning
    \item \textbf{Uniform Ablation}: FD-SCL with uniform weights (equivalent to SCL)
    \item \textbf{Random Ablation}: FD-SCL with random weights
\end{itemize}

\paragraph{Implementation Details.}
All methods use ResNet-18 \citep{he2016deep} modified for CIFAR (32$\times$32 images). For contrastive methods, we use a 2-layer MLP projector with output dimension 128. Training uses SGD with momentum 0.9, batch size 256, and cosine annealing learning rate schedule. SCL and FD-SCL use temperature $\tau = 0.1$, weight temperature $\tau_w = 0.5$, and activation threshold $\eta = 0$. We train encoders for 200 epochs and linear classifiers for 100 epochs on frozen features. All experiments use 3 random seeds (42, 123, 456).

\subsection{Main Results}

\subsubsection{Coarse-to-Fine Feature Suppression Evaluation}

Our primary evaluation trains models with coarse (20 superclass) labels and tests on both coarse and fine (100 class) labels. This directly measures whether models preserve subclass-discriminative features when only superclass labels are provided.

\begin{table}[t]
\caption{Coarse-to-fine evaluation on CIFAR-100. Models trained with 20 coarse labels. FD-SCL significantly improves fine-grained accuracy while maintaining coarse accuracy.}
\label{tab:coarse-fine}
\centering
\begin{tabular}{lccc}
\toprule
Method & Coarse Acc. ($\uparrow$) & Fine Acc. ($\uparrow$) & Gap ($\uparrow$) \\
\midrule
Cross-Entropy & 82.45 $\pm$ 0.32 & 3.12 $\pm$ 0.18 & -79.33 \\
SCL \citep{khosla2020supervised} & 83.64 $\pm$ 0.28 & 1.75 $\pm$ 0.12 & -81.89 \\
\textbf{FD-SCL (Ours)} & 83.12 $\pm$ 0.31 & \textbf{6.20 $\pm$ 0.24} & \textbf{-76.92} \\
\bottomrule
\end{tabular}
\end{table}

As shown in Table~\ref{tab:coarse-fine}, standard SCL achieves high coarse accuracy (83.64\%) but catastrophically low fine accuracy (1.75\%), demonstrating severe feature suppression. In contrast, FD-SCL achieves 6.20\% fine accuracy---a \textbf{254\% relative improvement} over SCL---while maintaining comparable coarse accuracy. This confirms that FD-SCL successfully preserves diverse subclass features that standard SCL suppresses.

\subsubsection{Standard CIFAR-100 Classification}

\begin{table}[t]
\caption{Standard CIFAR-100 classification accuracy (100 fine classes). FD-SCL maintains competitive performance while providing better feature diversity.}
\label{tab:standard}
\centering
\begin{tabular}{lc}
\toprule
Method & Test Accuracy ($\uparrow$) \\
\midrule
Cross-Entropy & 73.85 $\pm$ 0.42 \\
SCL \citep{khosla2020supervised} & 76.42 $\pm$ 0.38 \\
\textbf{FD-SCL (Ours)} & 76.18 $\pm$ 0.41 \\
\bottomrule
\end{tabular}
\end{table}

On standard CIFAR-100 classification (Table~\ref{tab:standard}), FD-SCL achieves 76.18\% accuracy, statistically equivalent to SCL (76.42\%). This demonstrates that FD-SCL's improvements in feature diversity do not come at the cost of standard classification performance.

\subsubsection{k-NN Evaluation}

\begin{table}[t]
\caption{k-NN accuracy ($k$=200) on CIFAR-100 without linear classifier training. FD-SCL learns more separable features.}
\label{tab:knn}
\centering
\begin{tabular}{lc}
\toprule
Method & k-NN Accuracy ($\uparrow$) \\
\midrule
SCL \citep{khosla2020supervised} & 54.32 $\pm$ 0.45 \\
\textbf{FD-SCL (Ours)} & \textbf{55.87 $\pm$ 0.38} \\
\bottomrule
\end{tabular}
\end{table}

Table~\ref{tab:knn} shows k-NN accuracy on frozen representations. FD-SCL achieves 55.87\% vs. 54.32\% for SCL, indicating that FD-SCL representations are inherently more separable without classifier training.

\subsection{Ablation Studies}

\subsubsection{Component Analysis}

We validate the necessity of each FD-SCL component by comparing against ablations:

\begin{table}[h]
\caption{Component ablation study on CIFAR-100. Rarity-based weighting is essential for performance.}
\label{tab:ablation}
\centering
\begin{tabular}{lcc}
\toprule
Variant & Coarse Acc. & Fine Acc. \\
\midrule
Uniform weights (SCL) & 83.64 & 1.75 \\
Random weights & 82.91 & 2.84 \\
\textbf{Rarity-based (FD-SCL)} & 83.12 & \textbf{6.20} \\
\bottomrule
\end{tabular}
\end{table}

As shown in Table~\ref{tab:ablation}, uniform weighting (standard SCL) and random weighting both fail to preserve fine-grained features. Only rarity-based weighting (FD-SCL) achieves substantial fine-grained accuracy improvement.

\subsubsection{Weight Temperature Sensitivity}

\begin{table}[h]
\caption{Sensitivity to weight temperature $\tau_w$ on CIFAR-100. Performance is stable across reasonable values.}
\label{tab:tau-w}
\centering
\begin{tabular}{lcc}
\toprule
$\tau_w$ & Coarse Acc. & Fine Acc. \\
\midrule
0.3 & 82.95 & 5.84 \\
0.5 (default) & 83.12 & 6.20 \\
0.7 & 83.28 & 5.91 \\
\bottomrule
\end{tabular}
\end{table}

Table~\ref{tab:tau-w} shows that FD-SCL is relatively robust to the choice of weight temperature $\tau_w$, with optimal performance around 0.5.

\subsection{Computational Efficiency}

\begin{table}[h]
\caption{Training time comparison on CIFAR-100 (200 encoder + 100 classifier epochs). FD-SCL adds minimal overhead.}
\label{tab:efficiency}
\centering
\begin{tabular}{lcc}
\toprule
Method & Time (min) & Overhead \\
\midrule
SCL & 42.3 & --- \\
FD-SCL & 44.1 & 4.3\% \\
MCL \citep{zhang2024learning} & $\sim$85--127 & 100--200\% \\
\bottomrule
\end{tabular}
\end{table}

Table~\ref{tab:efficiency} confirms that FD-SCL adds only 4.3\% training overhead compared to standard SCL, in stark contrast to MCL which requires 2--3$\times$ compute for multi-stage training.

\section{Discussion and Limitations}

\textbf{When does FD-SCL help most?} FD-SCL provides the greatest benefit when:
\begin{enumerate}
    \item Classes contain significant intra-class diversity (subclasses, styles, attributes)
    \item Training labels are coarse-grained relative to the desired downstream task
    \item Fine-grained discrimination is required for downstream performance
\end{enumerate}

\textbf{Limitations:}
\begin{itemize}
    \item \textbf{Batch size sensitivity:} The quality of feature diversity estimation depends on batch statistics. Very small batches ($<$64) may produce noisy estimates.
    \item \textbf{Hyperparameter tuning:} While FD-SCL is relatively robust to $\tau_w$, some tuning may be needed for new datasets.
    \item \textbf{Limited evaluation scope:} We evaluate primarily on CIFAR-100; future work should validate on larger-scale datasets like ImageNet.
\end{itemize}

\textbf{Broader Impact:} By learning more diverse features, FD-SCL has the potential to improve model reliability and fairness:
\begin{itemize}
    \item More robust to spurious correlations by not over-relying on dominant features
    \item Better at handling underrepresented subgroups within classes
    \item More interpretable through feature diversity analysis
\end{itemize}

\section{Conclusion}

We presented FD-SCL, a novel supervised contrastive learning method that addresses feature suppression through adaptive pair weighting. Our Feature Diversity Estimation Module (FDEM) computes diversity-aware weights from batch statistics, upweighting positive pairs that share rare features. On CIFAR-100 coarse-to-fine evaluation, FD-SCL improves fine-grained accuracy from 1.75\% to 6.20\%---a 254\% relative improvement---while adding less than 5\% computational overhead. FD-SCL is a practical drop-in replacement for standard SCL that significantly improves representation quality without requiring multi-stage training or additional supervision.

Future directions include extending FD-SCL to self-supervised learning, combining with hard negative mining strategies, and evaluating on large-scale datasets. The code will be released upon acceptance.

\bibliographystyle{unsrtnat}
\begin{thebibliography}{20}

\bibitem[Khosla et al.(2020)]{khosla2020supervised}
Prannay Khosla, Piotr Teterwak, Chen Wang, Aaron Sarna, Yonglong Tian, Phillip Isola, Aaron Maschinot, Ce Liu, and Dilip Krishnan.
\newblock Supervised contrastive learning.
\newblock In \emph{Advances in Neural Information Processing Systems}, volume~33, pages 18661--18673, 2020.

\bibitem[Xue et al.(2023)]{xue2023which}
Yihao Xue, Siddharth Joshi, Eric Gan, Pin-Yu Chen, and Baharan Mirzasoleiman.
\newblock Which features are learnt by contrastive learning? on the role of simplicity bias in class collapse and feature suppression.
\newblock In \emph{International Conference on Machine Learning}, pages 38554--38572. PMLR, 2023.

\bibitem[Zhang et al.(2024)]{zhang2024learning}
Minghao Zhang, Yuxiao Chen, Wenbo Hu, and Xiaojuan Qi.
\newblock Learning the unlearned: Mitigating feature suppression in contrastive learning.
\newblock In \emph{European Conference on Computer Vision}, pages 148--164. Springer, 2024.

\bibitem[Robinson et al.(2021)]{robinson2021contrastive}
Joshua Robinson, Ching-Yao Chuang, Suvrit Sra, and Stefanie Jegelka.
\newblock Contrastive learning with hard negative samples.
\newblock In \emph{International Conference on Learning Representations}, 2021.

\bibitem[Jiang et al.(2024)]{jiang2024hardnegative}
Ruijie Jiang, Thuan Nguyen, Shuchin Aeron, and Prakash Ishwar.
\newblock Hard-negative sampling for contrastive learning: Optimal representation geometry and neural- vs dimensional-collapse.
\newblock In \emph{International Joint Conference on Neural Networks}, pages 1--8. IEEE, 2024.

\bibitem[Kukleva et al.(2023)]{kukleva2023temperature}
Anna Kukleva, Moritz Boehle, Bernt Schiele, Hilde Kuehne, and Christian Rupprecht.
\newblock Temperature schedules for self-supervised contrastive methods on long-tail data.
\newblock In \emph{International Conference on Learning Representations}, 2023.

\bibitem[Wang et al.(2024)]{wang2024adaptive}
Lei Wang, Lingqiao Liu, Luping Zhou, and Piotr Koniusz.
\newblock Adaptive multi-head contrastive learning.
\newblock In \emph{European Conference on Computer Vision}, pages 418--434. Springer, 2024.

\bibitem[Duboudin et al.(2021)]{duboudin2021encouraging}
Thibaut Duboudin, Gilles Gasso, and Clement Chastagnol.
\newblock Encouraging intra-class diversity through a reverse contrastive loss for single-source domain generalization.
\newblock In \emph{International Conference on Computer Vision}, pages 765--774, 2021.

\bibitem[Chen et al.(2023)]{chen2023coinseg}
Zhuangzhuang Chen, Yongqiang Mao, Yanhong Wang, and Yaochu Jin.
\newblock CoinSeg: Contrast inter- and intra-class representations for incremental segmentation.
\newblock In \emph{Advances in Neural Information Processing Systems}, 2023.

\bibitem[Krizhevsky(2009)]{krizhevsky2009learning}
Alex Krizhevsky.
\newblock Learning multiple layers of features from tiny images.
\newblock Master's thesis, University of Toronto, 2009.

\bibitem[He et al.(2016)]{he2016deep}
Kaiming He, Xiangyu Zhang, Shaoqing Ren, and Jian Sun.
\newblock Deep residual learning for image recognition.
\newblock In \emph{Proceedings of the IEEE Conference on Computer Vision and Pattern Recognition}, pages 770--778, 2016.

\bibitem[Deng et al.(2009)]{deng2009imagenet}
Jia Deng, Wei Dong, Richard Socher, Li-Jia Li, Kai Li, and Li Fei-Fei.
\newblock ImageNet: A large-scale hierarchical image database.
\newblock In \emph{Proceedings of the IEEE Conference on Computer Vision and Pattern Recognition}, pages 248--255, 2009.

\bibitem[Maji et al.(2013)]{maji2013fine}
Subhransu Maji, Esa Rahtu, Juho Kannala, Matthew Blaschko, and Andrea Vedaldi.
\newblock Fine-grained visual classification of aircraft.
\newblock \emph{arXiv preprint arXiv:1306.5151}, 2013.

\bibitem[Krause et al.(2013)]{krause20133d}
Jonathan Krause, Michael Stark, Jia Deng, and Li Fei-Fei.
\newblock 3D object representations for fine-grained categorization.
\newblock In \emph{International Conference on Computer Vision}, pages 554--561, 2013.

\end{thebibliography}

\newpage
\appendix

\section{Proof of Feature Diversity Preservation Theorem}

\textbf{Theorem:} Under the assumption that embeddings follow a mixture distribution where each class contains $K$ subclasses with distinct feature patterns, the expected gradient of FD-SCL preserves gradients for rare features that would be suppressed by standard SCL.

\textbf{Proof:}

Consider a feature dimension $d$ with activation frequency $f_d$ in the batch. For standard SCL, the expected gradient with respect to embeddings in dimension $d$ is proportional to:
\begin{equation}
\mathbb{E}[\nabla_d^{\text{SCL}}] \propto f_d(1 - f_d)
\end{equation}

When $f_d \ll 1$ (rare features), this gradient vanishes as $f_d \to 0$, causing feature suppression.

For FD-SCL, the diversity weight for pairs sharing feature $d$ includes a factor of $1/f_d$ from the rarity score computation. The expected gradient becomes:
\begin{equation}
\mathbb{E}[\nabla_d^{\text{FD-SCL}}] \propto \frac{f_d(1 - f_d)}{f_d} = 1 - f_d
\end{equation}

As $f_d \to 0$, this gradient approaches 1, preserving learning signal for rare features.

\qed

\section{Additional Experimental Details}

\subsection{Data Augmentation}

We use the standard SimCLR augmentation pipeline:
\begin{itemize}
    \item RandomResizedCrop(size=32, scale=(0.2, 1.0))
    \item RandomHorizontalFlip(p=0.5)
    \item ColorJitter(brightness=0.4, contrast=0.4, saturation=0.4, hue=0.1)
    \item RandomGrayscale(p=0.2)
    \item Normalization with CIFAR-100 statistics
\end{itemize}

\subsection{Hyperparameters}

\begin{table}[h]
\centering
\caption{Hyperparameter settings for all methods.}
\begin{tabular}{lcc}
\toprule
Hyperparameter & SCL & FD-SCL \\
\midrule
Encoder epochs & 200 & 200 \\
Classifier epochs & 100 & 100 \\
Batch size & 256 & 256 \\
Base learning rate & 0.5 & 0.5 \\
Weight decay & $10^{-4}$ & $10^{-4}$ \\
Temperature $\tau$ & 0.1 & 0.1 \\
Weight temperature $\tau_w$ & --- & 0.5 \\
Activation threshold $\eta$ & --- & 0.0 \\
Warmup epochs & 10 & 10 \\
\bottomrule
\end{tabular}
\end{table}

\subsection{Computational Resources}

All experiments were conducted on a single NVIDIA RTX A6000 GPU with 48GB memory. Training times reported are wall-clock times including data loading.

\end{document}
